\subsection*{Part (c). Keeping the pretrained features}

\lead{Setup.} All runs start from the related-DNA checkpoint and train 16 passes on 30k chunks (Table~\ref{tab:p4c}). Variants: full fine-tuning at the scratch LR with WD 3.2 or .1; WD toward the initialization instead of toward zero (AdamW decay replaced by a pull to the pretrained weights); frozen head (train the body, keep embeddings and output layer); frozen body (train only embeddings, output layer and norms); and the same variants at LR $3\times10^{-4}$, WD 0, queued second and marked where unfinished.
\begin{itemize}
\item The pretrained body is the asset. Freezing it and training only the head reaches 1.2431, 0.02 below the best scratch model (1.2627) and below the supplied five-member ensemble (1.2543). Decaying toward the initialization is nearly as good (1.2452). Full fine-tuning at the same LR gives 1.2653, so the recipe that is right for scratch costs 0.02 nats when applied to a good start.
\item Frozen head (1.2624) is no better than full fine-tuning: the damage is done in the body, not the output layer, and training the body freely at LR .003 is what forgets.
\item Weak decay is worse still: WD .1 ends at 1.4431, the same memorization failure as from scratch. Pretraining does not protect against overfitting a repeated 30k-chunk pool.
\item Lowering the LR to .001 does not rescue full fine-tuning: that run reached 1.2618 at pass 8 and then drifted back to 1.2652 by pass 16, the same final value as LR .003. The forgetting is driven by the number of free updates on the repeated pool, not the step size alone: at LR $3\times10^{-4}$ and $10^{-4}$ with WD 0, 16 passes end at 1.64 and 1.38 after passing through 1.23 at pass 1. Decaying toward the initialization at $3\times10^{-4}$ was at 1.2309 after two passes when the queue was stopped, matching the one-pass optimum, which is the behaviour the constraint is meant to give.
\end{itemize}
